\chapter{熱力学}
\section{背景：保存される量と、変換できる量を分ける}
発電機関では燃料から得たエネルギーの全てを仕事へ変換できない。断熱材ではエネルギーが保存されても流れる速さが重要であり、圧力計ではその熱移動を信号として利用する。熱力学は、巨視的な系を少数の状態変数で表して、エネルギーの保存、変化の方向、平衡条件を調べる枠組みである。輸送現象はそこへ時間・空間の勾配を加える。

物理専攻の読者には、公式の再暗記よりも\emph{系の境界、仕事の向き、モデルの条件}を問題ごとに確認することが重要である。温度の数値が°Cであっても、状態方程式の比には絶対温度が必要である。「断熱」と「仕事をしない」、「理想気体」と「実在気体」、「比熱」と「熱伝導率」は異なる概念である。以下の理論から個別問題の条件へ接続する。

\section{第一法則：内部エネルギーと熱・仕事の符号}
\label{theory:thermo:firstlaw}
閉じた系について、系に入る熱量を$Q>0$、系へ外部から行われる仕事を$W_{\mathrm{on}}>0$とすると、全体の運動・位置エネルギー変化を無視した第一法則は
\[
 \Delta U=Q+W_{\mathrm{on}}=Q-W_{\mathrm{out}}
\]
である。$U$は状態量であり、$Q,W$は経路に依存する。体積仕事だけなら
\[
 \delta W_{\mathrm{on}}=-p_{\mathrm{ext}}\,dV,
 \qquad W_{\mathrm{out}}=\int p_{\mathrm{ext}}\,dV.
\]
圧縮では$dV<0$なので系へされる仕事は正である。準静的な力学平衡では外圧と系の圧力がほぼ一致するが、急速な膨張で系内の圧力をそのまま仕事の式へ代入してよいわけではない。符号と経路の扱いは大学講義の第一法則の整理と整合する\parencite{mit_thermo_firstlaw2008}。

断熱は境界からの熱流がない条件であり、電気入力や攪拌仕事を否定しない。断熱シリンダーの物質Aをヒーターで気化させる問題では、電気的な入力が物質Aのエネルギーを増やす。求める仕事が$p\Delta V$なら、ヒーターの総入力や潜熱と取り違えない。

\section{理想気体：状態方程式と熱容量の差}
\label{theory:thermo:ideal}
理想気体は分子間相互作用と占有体積の効果を無視するモデルである。量$n$を固定すると
\[
 pV=nRT,
 \qquad T[\mathrm{K}]=t[{}^\circ\mathrm{C}]+273.15.
\]
等温では$V\propto1/p$、定圧では$V\propto T$、定積では$p\propto T$となる。273へ丸めると本章の27°C、327°C、−123°Cは300、600、150Kとなり、候補の精度には十分である。温度差1°Cと1Kは等しいが、温度の比は°Cの値から作れない。

\subsection{定圧と定積で必要熱量が異なる理由}
\label{theory:thermo:capacity}
熱容量は系全体の温度変化当たりのエネルギー、比熱は質量当たり、モル熱容量は物質量当たりである。単相・一定組成の理想気体では$U=U(T)$であり、
\[
 C_{V,m}=\frac1n\left(\frac{\partial U}{\partial T}\right)_V,
 \qquad H=U+pV=U+nRT,
 \qquad C_{p,m}-C_{V,m}=R.
\]
定積加熱は内部エネルギーを増やし、定圧加熱はそれに加え膨張仕事を行う。このため同じ理想気体の同じ昇温に必要な熱は定圧の方が多い。熱容量がほぼ一定の区間なら$Q_V=nC_{V,m}\Delta T$、$Q_p=nC_{p,m}\Delta T$であり、大きく変わる区間では積分する。

\section{分子自由度と比熱：温度尺度を読む}
\label{theory:thermo:degrees}
分子のエネルギーには並進、回転、振動と、分子間の相互作用が寄与する。古典的な二次形式の自由度は1つ当たり平均$k_BT/2$を持つ。単原子理想気体では並進の3自由度から$C_{V,m}=3R/2$となる。二原子分子では回転の2自由度が励起され、振動が凍結している範囲なら
\[
 C_{V,m}=\frac52R,\qquad C_{p,m}=\frac72R,
 \qquad \gamma=\frac{C_{p,m}}{C_{V,m}}=\frac75.
\]
振動は運動とポテンシャルの双方を持つので、1振動モードが古典的に励起されると$R$の寄与を加える。低温では励起準位の間隔が$k_BT$より大きくなり、その自由度の熱容量への寄与は凍結する。従って「比較的低い温度で7/5」は、振動の励起温度より低く、回転の励起温度より高い中間範囲を読む必要がある。内部エネルギーの絶対零点を分子種間でそろえる問題と、温度によって増える熱エネルギーの自由度比較は区別する。

\section{断熱膨張：外部仕事がある場合と自由膨張}
\label{theory:thermo:adiabatic}
体積仕事だけの断熱過程では$nC_{V,m}dT=-p_{\mathrm{ext}}dV$である。外部へ正の仕事を行えば内部エネルギーが減り、理想気体は冷える。可逆で熱容量一定なら$p_{\mathrm{ext}}\simeq p=nRT/V$より
\[
 \frac{dT}{T}=-\frac{R}{C_{V,m}}\frac{dV}{V},
 \qquad TV^{\gamma-1}=\text{一定},\qquad pV^\gamma=\text{一定}.
\]
可逆断熱と、一定外圧に抗する不可逆断熱は区別して扱う\parencite{mit_thermo_adiabatic2008}。真空への自由膨張では外圧が0で仕事も0となり、断熱理想気体の$\Delta U$と$\Delta T$は0である。膨張が起きることだけから、常に温度低下すると推論してはいけない。

\begin{tcolorbox}[title=原典の限定不足を保持する]
平成25年度問9と平成30年度問6は、断熱膨張の記述に仕事又は可逆性の指定がない。定積・定圧の必要熱量を逆転した肢は明確に誤るが、断熱膨張の肢も自由膨張を含めれば無条件には正しくない。教材側で原典を可逆膨張へ書き換えず、独立判定と解釈疑義を保存する。
\end{tcolorbox}

\section{第二法則とカルノー機関：温度比の意味}
\label{theory:thermo:entropy}
熱が高温物体から低温物体へ受動的に流れるとき、エネルギーは保存されるが全体のエントロピーは増える。高温源$T_H$から熱$Q>0$が低温源$T_L$へ移れば
\[
 \Delta S_{\mathrm{total}}=-\frac{Q}{T_H}+\frac{Q}{T_L}>0.
\]
冷凍機は外部仕事で低温側の熱を運び出し、装置と周囲を含めた全体で第二法則を満たす。「低温から高温へ熱を移す装置がある」ことは受動的熱移動の法則への反例ではない。

\subsection{可逆サイクルの効率の上限}
\label{theory:thermo:carnot}
2熱源の間で作動する可逆サイクルでは、正の大きさを使って吸熱$Q_H$、放熱$Q_L$を表すと、
\[
 \frac{Q_H}{T_H}=\frac{Q_L}{T_L},
 \qquad W=Q_H-Q_L,
 \qquad \eta=\frac{W}{Q_H}=1-\frac{T_L}{T_H}.
\]
これがカルノー効率である。不可逆なサイクルでは、熱源も含めたエントロピー生成が非負であるため$Q_L/T_L-Q_H/T_H\ge0$、従って効率はこの値を超えない。$T_L\to T_H$で効率は0へ向かう。この極限は分母を温度の和にした誤式を見分ける手掛かりとなる。放熱の問題は$Q_L=Q_HT_L/T_H$、仕事の問題はその差を使う。

\section{仕事の経路と断熱混合の熱収支}
\label{theory:thermo:workpaths}
定圧・定積の段階を表へ分けると、各段階で不変な量を取り違えずに済む。固定$n$の定圧段階なら
\[
 W_{\mathrm{on}}=-p\Delta V=-nR\Delta T,
 \qquad nR=\frac{p_0V_0}{T_0}.
\]
定積段階の体積仕事は0である。これは「第3段階」の仕事を聞く問題で、サイクル全体の仕事や、初期圧力を全段階へ使う誤りを避ける。27°C、227°C、27°C、−73°Cの順の令和7年度問6では、第3段階の冷却幅は約100Kであり、初期量から$nR$を求めるだけでも仕事が出る。

気化の定圧仕事では、液体体積$V_l$を無視して理想蒸気の終体積を使えば
\[
 W_{\mathrm{out}}=p(V_g-V_l)\simeq pV_g=nRT.
\]
外圧が打ち消されるのは理想蒸気と体積無視という仮定の結果である。一般の相変化で圧力を無視してよいわけではない。

\subsection{熱容量を重みにする接触熱収支}
\label{theory:thermo:mixheat}
2物体を接触させ、相変化と外界への熱流を無視し比熱を一定とすると、
\[
 m_1c_1(T_f-T_1)+m_2c_2(T_f-T_2)=0,
 \qquad T_f=\frac{m_1c_1T_1+m_2c_2T_2}{m_1c_1+m_2c_2}.
\]
最終温度は熱容量を重みにした平均である。等質量でも比熱が違えば単純平均ではない。平成26年度問6は水の熱容量が金属球の10倍なので、最終温度は水の初期25°Cに近い。容器も吸熱する現実の実験では、その熱容量を収支へ加える。

\section{相平衡、混合、実在気体}
\label{theory:thermo:phase}
相平衡では各相の化学ポテンシャルが等しい。純物質で$d\mu=-s_m\,dT+v_m\,dp$を液相と気相へ適用し、共存曲線上で両者を等置すると
\[
 \frac{dp}{dT}=\frac{s_{g,m}-s_{l,m}}{v_{g,m}-v_{l,m}}
 =\frac{L_m}{T\Delta v_m}.
\]
これがクラペイロンの式である。液体の体積を無視し、気体を理想気体とすると$d\ln p/dT=L_m/(RT^2)$となる。さらに潜熱$L_m$を区間内で一定と近似すれば
\[
 \ln\frac{p_2}{p_1}=-\frac{L_m}{R}\left(\frac1{T_2}-\frac1{T_1}\right).
\]
このクラウジウス＝クラペイロンの関係で外圧から沸点を求める\parencite{mit_clapeyron2008}。外圧が低い山頂では沸点が低い。三重点の温度・圧力だけで共存曲線の全体が一意に決まるわけではなく、潜熱などの情報が必要である。式の名前を選ぶ問題へ、指定されない山頂の沸点の数値を補わない。

一次相転移は通常、$G$の一階微分である$S,V$などに不連続性を持ち、潜熱は$T\Delta S$である。二次・連続転移では一階微分が連続でも二階微分に跳びや発散が現れ得る。従って二次転移なら比熱も必ず連続という主張は誤る。過冷却は核生成が遅れて準安定相が続く速度論的な状態であり、平衡相図と区別する。

\subsection{異種気体の混合と理想溶液}
\label{theory:thermo:mixing}
同温同圧の異種理想気体を混ぜると、組成$x_i$について$\Delta S_{\mathrm{mix}}=-nR\sum_i x_i\ln x_i>0$、$\Delta G_{\mathrm{mix}}=nRT\sum_i x_i\ln x_i<0$となる。理想溶液では液相の活動度がモル分率と等しくなり、蒸気圧は$p_i=x_i p_i^*$に従う。これはラウールの法則であり、現実の分子間相互作用が大きく異なる溶液では活動度の補正が必要になる\parencite{mit_solutions2008}。同じ物質の分割を混ぜることを異種混合と数えることはしない。

\subsection{実在気体の圧縮係数}
\label{theory:thermo:realgas}
実在気体の圧縮係数$Z=pV_m/(RT)$は、同じ$T,p$の理想モル体積に対する実在モル体積の比である。ファンデルワールス式なら
\[
 \left(p+\frac{a}{V_m^2}\right)(V_m-b)=RT,
 \qquad Z=\frac{V_m}{V_m-b}-\frac{a}{RTV_m}.
\]
第1項の排除体積は$Z$を増やし、第2項の引力は減らす。温度・密度によって$Z$は1より小さくも大きくもなり、希薄極限で1へ向かう。この式は定性的にずれを理解する近似であり、実用の高精度状態方程式を代替する保証ではない。

\section{熱伝導・放射・計測：輸送の担体を区別する}
\label{theory:thermo:transport}
線形・等方の熱伝導では熱流密度$\boldsymbol{j}_q=-\kappa\nabla T$である。熱伝導率$\kappa$の単位はW/(m・K)、比熱の単位はJ/(kg・K)であり、前者は輸送のしやすさ、後者は温度を変えるエネルギーを表す。結晶では格子振動が熱を運び、金属では伝導電子も重要である。通常の縮退金属の電子輸送では
\[
 \frac{\kappa_e}{\sigma T}\simeq L_0,
 \qquad L_0=\frac{\pi^2}{3}\left(\frac{k_B}{e}\right)^2
\]
というウィーデマン＝フランツ則が成り立つ近似がある。格子伝導が大きい材料まで、電気伝導の大小から熱伝導を判定できるわけではない。

産総研の室温付近の値では、Alは226〜237、Feは72〜80.4、Cuは386〜402、Agは429W/(m・K)程度である\parencite{aist_metal_thermal}。高品質バルクダイヤモンドは格子による熱伝導が大きく、通常の銅・鉄鋼・板ガラスを上回るが、粒径や欠陥で変わる\parencite{wang_sun_diamond2024}。板ガラスの目安は約1W/(m・K)である\parencite{agc_glass_thermal}。これらは試料の代表値であり、物質名だけで全ての試料の厳密値を決める表ではない。

多孔質の断熱材は固体の伝導経路を減らし、気体の小さい熱伝導を利用する。ただし孔が大きく対流が起きる場合や高温での放射は別に評価する。熱放射は電磁波によるエネルギー移動であり、物質媒体の移動を必要としない。

\subsection{熱輸送と熱電能を利用する計測}
\label{theory:thermo:measurement}
ピラニ真空計は加熱フィラメントから気体へ逃げる熱が圧力で変わることを利用する。温度維持の加熱電力又は抵抗変化を圧力へ換算する原理である\parencite{ulvac_pirani}。吸着で生じる圧力変化から熱伝導率を測る説明とは、入力と出力が違う。実測では気体種や校正範囲も確認する。

熱電対は異種材料の熱電能の差によるゼーベック効果を用いる。基準温度$T_r$と測定温度$T_h$に対する電圧は、向きの定義に応じて$\int_{T_r}^{T_h}(S_A-S_B)\,dT$である。熱起電力を読むには基準側の温度も必要である。電流を流した接合で熱が出入りするペルティエ効果とは役割を区別する。

\section{固体比熱：電子・格子・相転移を分ける}
\label{theory:thermo:solidcapacity}
電子の励起はフェルミ準位の近傍に限られ、滑らかな状態密度$D(E_F)$を持つ通常の金属の低温電子比熱は
\[
 C_e=\gamma T,\qquad\gamma=\frac{\pi^2}{3}k_B^2D(E_F).
\]
$D$を系全体の状態数/エネルギーと定義すれば$C_e$も系全体の熱容量である。単位体積当たりの状態密度なら体積当たりとなる。反比例ではなく比例であり、熱的に励起できる状態が多いほど比熱が大きい\parencite{mit_electronic_heat2012}。

格子ではデバイ温度$\Theta_D$が振動の励起の尺度を与え、1molの原子の比熱はデバイモデルで
\[
 C_{V,m}=9R\left(\frac{T}{\Theta_D}\right)^3
 \int_0^{\Theta_D/T}\frac{x^4 e^x}{(e^x-1)^2}\,dx.
\]
低温で$T^3$、高温で$3R$へ近づく。このデュロン＝プティ則の高温とは、室温という固定値でなく$T\gg\Theta_D$の意味である\parencite{mit_debye2005}。ダイヤモンドのように励起温度の高い物質では室温の格子比熱は$3R$からずれる。格子比熱が小さいことと熱伝導率が大きいことは矛盾しない。後者には熱を運ぶ励起の速度と散乱の少なさも効く。

固体の熱容量差は、体積膨張率$\alpha_V=(1/V)(\partial V/\partial T)_p$、等温圧縮率$\kappa_T=-(1/V)(\partial V/\partial p)_T$を使って
\[
 C_p-C_V=\frac{TV\alpha_V^2}{\kappa_T}
\]
と表せる。一般に低温で熱膨張が小さくなるため差は小さくなるが、「100K以下なら全物質で無視できる」という普遍的境界ではない。転移近傍や特殊な材料の例外を無視しない。

\section{問題の読み方とこの初稿の確認範囲}
計算問題では、対象量を仕事・熱・内部エネルギーに分け、向きと単位を決めてから不変量を使う。知識問題では、理想モデル、温度範囲、相転移の次数、計測の入出力を確認する。複数の近似を含む5肢の問題で、最も明確な誤りの選択と他肢の普遍的な正しさは同じ判断ではない。令和3年度問7の温度範囲の省略は、断熱膨張の2問とともに査読事項として保持する。

このバッチは熱力学22問を原典画像から確認した初稿である。頻度表は確認済み部分集合と全体の暫定集計を区別し、将来出題を保証しない。以下の個別解説は全てAstra査読前であり、他章の未執筆問題も含めて作業を継続する。
\section{熱力学・熱機関}

\subsection{平成23年度 問5}
\label{q:h23-reg-applied-05}

\noindent\textbf{原典：} sources/H23/17\_applied\_science.pdf、PDF物理2頁。

\textbf{公式正答：} 1（single）。

\paragraph{問題の要旨} 気体の状態変化、第一法則、第二法則の5記述から最も不適切なものを選ぶ。

\Needspace{4\baselineskip}\paragraph{確認した条件・記述（独自要約）}\mbox{}\par

\noindent\textbf{方向：} 最も不適切な記述を選ぶ。\par

\noindent\textbf{原典の範囲：} 一定質量の気体、温度又は圧力一定。理想気体という語は明記されていない。\par

\paragraph{問われる核心} 第一法則の仕事の符号と気体の状態方程式

\paragraph{必要知識} 閉じた系、理想気体近似、絶対温度、熱と仕事の符号、受動的な熱移動。

\paragraph{考え方・正答を選ぶ理由（査読待ち初稿）} 等温の気体の体積と圧力の関係を\(pV=nRT\)で判定する。体積は圧力へ反比例するので1が不適切である。2は理想気体近似のシャルルの法則、3と4はエネルギー保存、5は第二法則である。気体法則は全ての実在気体へ厳密に成立するという主張へ拡張しない。

\Needspace{4\baselineskip}\paragraph{選択肢の検討（査読待ち初稿）}\mbox{}\par

\noindent\textbf{1：誤り} 温度一定・一定質量の気体の体積は圧力に比例する。理想気体なら\(pV=nRT\)であり、固定\(n,T\)では\(V=nRT/p\)と反比例する。比例という記述が誤りである。\par

\noindent\textbf{2：理想気体近似で正しい} 圧力一定・一定質量の気体の体積は絶対温度に比例する。理想気体又はその近似が成り立つ範囲では固定\(n,p\)により\(V=(nR/p)T\)である。\par

\noindent\textbf{3：正しい} 外部からの熱と系へされた仕事の和が内部エネルギー増加に等しい。熱を流入正、仕事を系へ入力正とすれば\(\Delta U=Q+W_{\mathrm{on}}\)である。系全体の運動・位置エネルギー変化はこの閉じた熱力学系では考えない。\par

\noindent\textbf{4：正しい} エネルギー変換で関与する全エネルギーの和が保存される。第一法則のエネルギー保存の内容である。対象系から外へ出るエネルギーを数え落とさない。\par

\noindent\textbf{5：正しい} 熱は高温から低温へ移り、自然に低温から高温へは移らない。第二法則の巨視的な受動的熱移動の方向である。冷凍機は外部仕事を使うので反例ではない。\par

\paragraph{再利用する考え方} モデルの前提を明示し、第一法則の仕事の符号と気体の状態方程式から記述や数値を判定する。

\paragraph{共通理論} \hyperref[theory:thermo:firstlaw]{第\ref*{theory:thermo:firstlaw}節}, \hyperref[theory:thermo:ideal]{第\ref*{theory:thermo:ideal}節}, \hyperref[theory:thermo:entropy]{第\ref*{theory:thermo:entropy}節}

\par 参考：\cite{mit_thermo_firstlaw2008}。

\paragraph{関連問題} \hyperref[q:h26-reg-applied-08]{H26-REG-APPLIED-08}

\par 難易度 1（校準初稿）、優先度未評価。

\subsection{平成24年度 問8}
\label{q:h24-reg-applied-08}

\noindent\textbf{原典：} sources/H24/17\_applied\_science.pdf、PDF物理2頁。

\textbf{公式正答：} 3（single）。

\paragraph{問題の要旨} 内部エネルギーの要素と単原子・二原子分子の自由度を、4つの空欄へ対応させる。

\Needspace{4\baselineskip}\paragraph{確認した条件・記述（独自要約）}\mbox{}\par

\noindent\textbf{A：} 分子間の力に由来するエネルギー。\par

\noindent\textbf{B：} 分子の熱運動によるエネルギー。\par

\noindent\textbf{C：} 二原子分子ではBに加わる分子内自由度。\par

\noindent\textbf{D：} 同温での分子1個当たりの二原子分子の平均内部エネルギーを、単原子分子と比較する。\par

\paragraph{問われる核心} 分子自由度と内部エネルギーの構成

\paragraph{必要知識} 並進・回転・振動、分子間相互作用、等分配と量子凍結。

\paragraph{考え方・正答を選ぶ理由（査読待ち初稿）} 分子間の力は位置エネルギー、分子の並進は運動エネルギーに対応する。二原子分子には回転と振動があり、励起される自由度は熱エネルギーに寄与するのでA位置・B運動・C振動回転・D大きいの3となる。量子論では低温に振動・回転が凍結し、常に全自由度が等分配されるわけではない。また分子種間の内部エネルギーの絶対零点を比較する問題とは区別する。

\Needspace{4\baselineskip}\paragraph{選択肢の検討（査読待ち初稿）}\mbox{}\par

\noindent\textbf{1：条件に不一致} A運動／B振動・回転／C位置／D小さい。Aは分子間相互作用の位置エネルギーであり、Bは並進の運動エネルギーである。A〜Cの対応が違う。\par

\noindent\textbf{2：条件に不一致} A運動／B振動・回転／C位置／D大きい。大きいというDだけでは足りず、A〜Cの対応が違う。\par

\noindent\textbf{3：条件に一致} A位置／B運動／C振動・回転／D大きい。A位置、B運動、C振動・回転、D大きいという対応に一致する。\par

\noindent\textbf{4：条件に不一致} A位置／B運動／C振動・回転／D小さい。二原子分子の回転等が熱的に励起される条件では平均の熱エネルギーが加わる。小さいは逆である。\par

\noindent\textbf{5：条件に不一致} A運動／B位置／C振動・回転／D大きい。相互作用の位置エネルギーと並進の運動エネルギーを逆にしている。\par

\paragraph{再利用する考え方} モデルの前提を明示し、分子自由度と内部エネルギーの構成から記述や数値を判定する。

\paragraph{共通理論} \hyperref[theory:thermo:degrees]{第\ref*{theory:thermo:degrees}節}

\par 参考：\cite{mit_debye2005}。

\par 難易度 1（校準初稿）、優先度未評価。

\subsection{平成25年度 問8}
\label{q:h25-reg-applied-08}

\noindent\textbf{原典：} sources/H25/17\_applied\_science.pdf、PDF物理4頁。

\textbf{公式正答：} 5（single）。

\paragraph{問題の要旨} 高温・低温の2熱源で作動する熱機関の最大効率を、機関名と式の組合せで選ぶ。

\Needspace{4\baselineskip}\paragraph{確認した条件・記述（独自要約）}\mbox{}\par

\noindent\textbf{温度：} 絶対温度\(T_H>T_L>0\)。\par

\noindent\textbf{効率：} 外部へ行う仕事／高温熱源から吸収する熱。\par

\noindent\textbf{空欄：} アは最大効率を与える熱機関名、イはその効率。\par

\paragraph{問われる核心} 第二法則が決める熱機関効率の上限

\paragraph{必要知識} 2熱源、可逆性、熱収支、絶対温度、サイクルのエントロピー収支。

\paragraph{考え方・正答を選ぶ理由（査読待ち初稿）} 熱収支は\(W=Q_H-Q_L\)、可逆サイクルでは\(Q_H/T_H=Q_L/T_L\)である。従って\(\eta=W/Q_H=1-T_L/T_H=(T_H-T_L)/T_H\)となり、カルノー機関とこの式を組み合わせた5である。不可逆損失がある実機の効率はこの上限以下になる。

\Needspace{4\baselineskip}\paragraph{選択肢の検討（査読待ち初稿）}\mbox{}\par

\noindent\textbf{1：条件に不一致} カルノー／\((T_H-T_L)/(T_H+T_L)\)。機関名は合うが効率の分母は吸熱源の絶対温度であり、温度の和ではない。\par

\noindent\textbf{2：条件に不一致} フェルマー／\((T_H-T_L)/(T_H+T_L)\)。機関名と分母の両方が違う。\par

\noindent\textbf{3：条件に不一致} カルノー／\(T_H/(T_H+T_L)\)。低温源が高温源に近づくと効率は0へ向かうべきであり、この式はその条件を満たさない。\par

\noindent\textbf{4：条件に不一致} フェルマー／\((T_H-T_L)/T_H\)。効率式は合うが機関名が違う。\par

\noindent\textbf{5：条件に一致} カルノー／\((T_H-T_L)/T_H\)。可逆なカルノー機関とその効率に一致する。\par

\paragraph{再利用する考え方} モデルの前提を明示し、第二法則が決める熱機関効率の上限から記述や数値を判定する。

\paragraph{共通理論} \hyperref[theory:thermo:carnot]{第\ref*{theory:thermo:carnot}節}, \hyperref[theory:thermo:entropy]{第\ref*{theory:thermo:entropy}節}

\paragraph{関連問題} \hyperref[q:h28-reg-applied-08]{H28-REG-APPLIED-08}

\par 難易度 1（校準初稿）、優先度未評価。

\subsection{平成26年度 問6}
\label{q:h26-reg-applied-06}

\noindent\textbf{原典：} sources/H26/17\_applied\_science.pdf、PDF物理3頁。

\textbf{公式正答：} 1（single）。

\paragraph{問題の要旨} 水と高温金属球を外界との熱交換なしで接触させ、共通の最終温度を求める。

\Needspace{4\baselineskip}\paragraph{確認した条件・記述（独自要約）}\mbox{}\par

\noindent\textbf{水：} 100g、25°C、比熱4.2J/(g・K)\par

\noindent\textbf{金属：} 100g、80°C、比熱0.42J/(g・K)\par

\noindent\textbf{無視するもの：} ビーカーへの伝熱と外界への熱損失。\par

\paragraph{問われる核心} 断熱混合のエネルギー保存と熱容量の重み

\paragraph{必要知識} 熱容量mc、比熱と熱容量の違い、相変化なし、熱収支。

\paragraph{考え方・正答を選ぶ理由（査読待ち初稿）} 水が受け取る熱と金属が放出する熱を等置し、\(100\times4.2(T-25)=100\times0.42(80-T)\)とする。整理すると\(420(T-25)=42(80-T)\)、\(462T=13860\)より\(T=30\,{}^\circ\mathrm{C}\)となる。第1候補である。相変化はなく温度差だけなので°Cの差をそのまま使える。

\Needspace{4\baselineskip}\paragraph{選択肢の検討（査読待ち初稿）}\mbox{}\par

\noindent\textbf{1：条件に一致} 約30°C。水420J/Kと金属42J/Kの熱容量で重みを付けると30°Cになる。\par

\noindent\textbf{2：条件に不一致} 約40°C。金属の熱容量は水の1/10であり、共通温度は水の初期温度にもっと近い。\par

\noindent\textbf{3：条件に不一致} 約50°C。等質量でも比熱が異なるため単純平均ではない。\par

\noindent\textbf{4：条件に不一致} 約60°C。水の大きな熱容量を無視している。\par

\noindent\textbf{5：条件に不一致} 約70°C。水の大きな熱容量を無視している。\par

\paragraph{再利用する考え方} モデルの前提を明示し、断熱混合のエネルギー保存と熱容量の重みから記述や数値を判定する。

\paragraph{共通理論} \hyperref[theory:thermo:mixheat]{第\ref*{theory:thermo:mixheat}節}

\par 難易度 2（校準初稿）、優先度未評価。

\subsection{平成26年度 問8}
\label{q:h26-reg-applied-08}

\noindent\textbf{原典：} sources/H26/17\_applied\_science.pdf、PDF物理3頁。

\textbf{公式正答：} 1（single）。

\paragraph{問題の要旨} 気体の状態変化、第一法則、第二法則の5記述から最も不適切なものを選ぶ。

\Needspace{4\baselineskip}\paragraph{確認した条件・記述（独自要約）}\mbox{}\par

\noindent\textbf{方向：} 最も不適切な記述を選ぶ。\par

\noindent\textbf{原典の範囲：} 一定質量の気体、温度又は圧力一定。理想気体という語は明記されていない。\par

\paragraph{問われる核心} 第一法則の仕事の符号と気体の状態方程式

\paragraph{必要知識} 閉じた系、理想気体近似、絶対温度、熱と仕事の符号、受動的な熱移動。

\paragraph{考え方・正答を選ぶ理由（査読待ち初稿）} 等温の気体の体積と圧力の関係を\(pV=nRT\)で判定する。体積は圧力へ反比例するので1が不適切である。2は理想気体近似のシャルルの法則、3と4はエネルギー保存、5は第二法則である。気体法則は全ての実在気体へ厳密に成立するという主張へ拡張しない。

\Needspace{4\baselineskip}\paragraph{選択肢の検討（査読待ち初稿）}\mbox{}\par

\noindent\textbf{1：誤り} 温度一定・一定質量の気体の体積は圧力に比例する。理想気体なら\(pV=nRT\)であり、固定\(n,T\)では\(V=nRT/p\)と反比例する。比例という記述が誤りである。\par

\noindent\textbf{2：理想気体近似で正しい} 圧力一定・一定質量の気体の体積は絶対温度に比例する。理想気体又はその近似が成り立つ範囲では固定\(n,p\)により\(V=(nR/p)T\)である。\par

\noindent\textbf{3：正しい} 外部からの熱と系へされた仕事の和が内部エネルギー増加に等しい。熱を流入正、仕事を系へ入力正とすれば\(\Delta U=Q+W_{\mathrm{on}}\)である。系全体の運動・位置エネルギー変化はこの閉じた熱力学系では考えない。\par

\noindent\textbf{4：正しい} エネルギー変換で関与する全エネルギーの和が保存される。第一法則のエネルギー保存の内容である。対象系から外へ出るエネルギーを数え落とさない。\par

\noindent\textbf{5：正しい} 熱は高温から低温へ移り、自然に低温から高温へは移らない。第二法則の巨視的な受動的熱移動の方向である。冷凍機は外部仕事を使うので反例ではない。\par

\paragraph{再利用する考え方} モデルの前提を明示し、第一法則の仕事の符号と気体の状態方程式から記述や数値を判定する。

\paragraph{共通理論} \hyperref[theory:thermo:firstlaw]{第\ref*{theory:thermo:firstlaw}節}, \hyperref[theory:thermo:ideal]{第\ref*{theory:thermo:ideal}節}, \hyperref[theory:thermo:entropy]{第\ref*{theory:thermo:entropy}節}

\par 参考：\cite{mit_thermo_firstlaw2008}。

\paragraph{関連問題} \hyperref[q:h23-reg-applied-05]{H23-REG-APPLIED-05}, \hyperref[q:r01-ret-applied-04]{R01-RET-APPLIED-04}

\par 難易度 1（校準初稿）、優先度未評価。

\subsection{平成27年度 問8}
\label{q:h27-reg-applied-08}

\noindent\textbf{原典：} sources/H27/17\_applied\_science.pdf、PDF物理3頁。

\textbf{公式正答：} 3（single）。

\paragraph{問題の要旨} カルノーサイクルの吸熱と2熱源の温度から、1サイクルの低温熱源への放熱を求める。

\Needspace{4\baselineskip}\paragraph{確認した条件・記述（独自要約）}\mbox{}\par

\noindent\textbf{高温熱源：} 1500K\par

\noindent\textbf{低温熱源：} 1200K\par

\noindent\textbf{高温側吸熱：} 750J\par

\noindent\textbf{求める量：} 低温源へ放出する熱量の正の大きさ。\par

\paragraph{問われる核心} カルノーサイクルの可逆エントロピー収支

\paragraph{必要知識} 絶対温度、吸熱と放熱の正の大きさ、仕事と熱収支。

\paragraph{考え方・正答を選ぶ理由（査読待ち初稿）} 可逆サイクルのエントロピー変化は0なので\(Q_L=Q_HT_L/T_H\)である。数値を代入すると\(750\times1200/1500=600\,\mathrm{J}\)となり、3を選ぶ。取り出した仕事は吸熱と放熱の差であり、求める放熱と混同しない。

\Needspace{4\baselineskip}\paragraph{選択肢の検討（査読待ち初稿）}\mbox{}\par

\noindent\textbf{1：条件に不一致} 480J。可逆性から得る低温側放熱600Jと一致しない。\par

\noindent\textbf{2：条件に不一致} 540J。可逆性から得る低温側放熱600Jと一致しない。\par

\noindent\textbf{3：条件に一致} 600J。吸熱に低温／高温の温度比を掛けた放熱に一致する。\par

\noindent\textbf{4：条件に不一致} 660J。可逆性から得る低温側放熱600Jと一致しない。\par

\noindent\textbf{5：条件に不一致} 720J。可逆性から得る低温側放熱600Jと一致しない。\par

\paragraph{再利用する考え方} モデルの前提を明示し、カルノーサイクルの可逆エントロピー収支から記述や数値を判定する。

\paragraph{共通理論} \hyperref[theory:thermo:carnot]{第\ref*{theory:thermo:carnot}節}

\paragraph{関連問題} \hyperref[q:r04-reg-applied-06]{R04-REG-APPLIED-06}

\par 難易度 2（校準初稿）、優先度未評価。

\subsection{平成28年度 問8}
\label{q:h28-reg-applied-08}

\noindent\textbf{原典：} sources/H28/17\_applied\_science.pdf、PDF物理3頁。

\textbf{公式正答：} 5（single）。

\paragraph{問題の要旨} 高温・低温の2熱源で作動する熱機関の最大効率を、機関名と式の組合せで選ぶ。

\Needspace{4\baselineskip}\paragraph{確認した条件・記述（独自要約）}\mbox{}\par

\noindent\textbf{温度：} 絶対温度\(T_H>T_L>0\)。\par

\noindent\textbf{効率：} 外部へ行う仕事／高温熱源から吸収する熱。\par

\noindent\textbf{空欄：} アは最大効率を与える熱機関名、イはその効率。\par

\paragraph{問われる核心} 第二法則が決める熱機関効率の上限

\paragraph{必要知識} 2熱源、可逆性、熱収支、絶対温度、サイクルのエントロピー収支。

\paragraph{考え方・正答を選ぶ理由（査読待ち初稿）} 熱収支は\(W=Q_H-Q_L\)、可逆サイクルでは\(Q_H/T_H=Q_L/T_L\)である。従って\(\eta=W/Q_H=1-T_L/T_H=(T_H-T_L)/T_H\)となり、カルノー機関とこの式を組み合わせた5である。不可逆損失がある実機の効率はこの上限以下になる。

\Needspace{4\baselineskip}\paragraph{選択肢の検討（査読待ち初稿）}\mbox{}\par

\noindent\textbf{1：条件に不一致} カルノー／\((T_H-T_L)/(T_H+T_L)\)。機関名は合うが効率の分母は吸熱源の絶対温度であり、温度の和ではない。\par

\noindent\textbf{2：条件に不一致} フェルマー／\((T_H-T_L)/(T_H+T_L)\)。機関名と分母の両方が違う。\par

\noindent\textbf{3：条件に不一致} カルノー／\(T_H/(T_H+T_L)\)。低温源が高温源に近づくと効率は0へ向かうべきであり、この式はその条件を満たさない。\par

\noindent\textbf{4：条件に不一致} フェルマー／\((T_H-T_L)/T_H\)。効率式は合うが機関名が違う。\par

\noindent\textbf{5：条件に一致} カルノー／\((T_H-T_L)/T_H\)。可逆なカルノー機関とその効率に一致する。\par

\paragraph{再利用する考え方} モデルの前提を明示し、第二法則が決める熱機関効率の上限から記述や数値を判定する。

\paragraph{共通理論} \hyperref[theory:thermo:carnot]{第\ref*{theory:thermo:carnot}節}, \hyperref[theory:thermo:entropy]{第\ref*{theory:thermo:entropy}節}

\paragraph{関連問題} \hyperref[q:h25-reg-applied-08]{H25-REG-APPLIED-08}

\par 難易度 1（校準初稿）、優先度未評価。

\subsection{平成29年度 問6}
\label{q:h29-reg-applied-06}

\noindent\textbf{原典：} sources/H29/17\_applied\_science.pdf、PDF物理2頁。

\textbf{公式正答：} 4（single）。

\paragraph{問題の要旨} ヒーターで1molの物質Aを一定外圧で全て気体にしたとき、物質Aが外部へ行う仕事を求める。

\Needspace{4\baselineskip}\paragraph{確認した条件・記述（独自要約）}\mbox{}\par

\noindent\textbf{系：} 断熱シリンダー内の1molの液体物質Aと、加熱用ヒーター。\par

\noindent\textbf{温度：} 物質Aの沸点373K\par

\noindent\textbf{外圧：} 0.1MPa一定\par

\noindent\textbf{仮定：} 液体体積を無視、気体状態の物質Aは理想気体。\par

\noindent\textbf{気体定数：} 8.3J/(mol・K)\par

\noindent\textbf{求める量：} 物質Aが外部へ行う仕事の正の大きさ。\par

\paragraph{問われる核心} 気化で増えた体積による定圧仕事

\paragraph{必要知識} 一定外圧の仕事、理想気体、液体体積無視、壁からの熱流と電気入力の区別。

\paragraph{考え方・正答を選ぶ理由（査読待ち初稿）} 外部への体積仕事は\(W_{\mathrm{out}}=p_{\mathrm{ext}}(V_g-V_l)\)である。液体体積を無視し、終状態の理想気体に\(p_{\mathrm{ext}}V_g=nRT\)を適用すると\(W_{\mathrm{out}}=1\times8.3\times373=3095.9\,\mathrm{J}\)である。最も近い3100Jの4を選ぶ。加熱の総エネルギーや潜熱を問う問題ではない。断熱壁でもヒーターから仕事としてエネルギーが入る。

\Needspace{4\baselineskip}\paragraph{選択肢の検討（査読待ち初稿）}\mbox{}\par

\noindent\textbf{1：条件に不一致} 0.023J。外圧×体積増加=nRTに対して桁又は値が合わない。\par

\noindent\textbf{2：条件に不一致} 830J。外圧×体積増加=nRTに対して桁又は値が合わない。\par

\noindent\textbf{3：条件に不一致} 2300J。外圧×体積増加=nRTに対して桁又は値が合わない。\par

\noindent\textbf{4：条件に一致} 3100J。理想蒸気のpV=\allowbreak{}nRTに一致する。\par

\noindent\textbf{5：条件に不一致} 8300J。外圧×体積増加=nRTに対して桁又は値が合わない。\par

\paragraph{再利用する考え方} モデルの前提を明示し、気化で増えた体積による定圧仕事から記述や数値を判定する。

\paragraph{共通理論} \hyperref[theory:thermo:firstlaw]{第\ref*{theory:thermo:firstlaw}節}, \hyperref[theory:thermo:workpaths]{第\ref*{theory:thermo:workpaths}節}

\par 参考：\cite{mit_thermo_firstlaw2008}。

\paragraph{関連問題} \hyperref[q:r06-reg-applied-06]{R06-REG-APPLIED-06}

\par 難易度 2（校準初稿）、優先度未評価。

\subsection{平成30年度 問7}
\label{q:h30-reg-applied-07}

\noindent\textbf{原典：} sources/H30/17\_applied\_science.pdf、PDF物理3頁。

\textbf{公式正答：} 1（single）。

\paragraph{問題の要旨} 室温での熱伝導率を4材料について大きい順に並べる。

\Needspace{4\baselineskip}\paragraph{確認した条件・記述（独自要約）}\mbox{}\par

\noindent\textbf{条件：} 室温で通常のバルク材料の代表的な熱伝導率を比較する。\par

\paragraph{問われる核心} 熱を運ぶ電子と格子振動・材料依存

\paragraph{必要知識} 熱伝導率と電気伝導率の区別、室温、バルク材料、純度と欠陥。

\paragraph{考え方・正答を選ぶ理由（査読待ち初稿）} ダイヤモンドの結晶格子は室温で熱をよく運び、高品質バルクでは銅より熱伝導率が大きい。通常の銅は鉄鋼より大きく、板ガラスは小さい。従ってダイヤモンド→銅→鉄鋼→板ガラスの1である。電気絶縁体なら熱伝導率も低い、という推論はダイヤモンドでは成立しない。結晶欠陥・粒径・純度・合金種で値は変わるので、全ての試料に同じ序列を断定しない。

\Needspace{4\baselineskip}\paragraph{選択肢の検討（査読待ち初稿）}\mbox{}\par

\noindent\textbf{1：条件に一致} ダイヤモンド→銅→鉄鋼→板ガラス。代表的な室温の大小関係に一致する。\par

\noindent\textbf{2：条件に不一致} ダイヤモンド→鉄鋼→銅→板ガラス。鉄鋼と銅が逆である。\par

\noindent\textbf{3：条件に不一致} 銅→鉄鋼→ダイヤモンド→板ガラス。高品質のバルクダイヤモンドを金属より下に置いている。\par

\noindent\textbf{4：条件に不一致} 銅→鉄鋼→板ガラス→ダイヤモンド。バルクダイヤモンドを最小にしている。\par

\noindent\textbf{5：条件に不一致} 銅→ダイヤモンド→鉄鋼→板ガラス。バルクダイヤモンドと銅が逆である。\par

\paragraph{再利用する考え方} モデルの前提を明示し、熱を運ぶ電子と格子振動・材料依存から記述や数値を判定する。

\paragraph{共通理論} \hyperref[theory:thermo:transport]{第\ref*{theory:thermo:transport}節}

\par 参考：\cite{aist_metal_thermal,wang_sun_diamond2024,agc_glass_thermal}。

\par 難易度 1（校準初稿）、優先度未評価。

\subsection{令和元年度 問6}
\label{q:r01-reg-applied-06}

\noindent\textbf{原典：} sources/R01/17\_applied\_science.pdf、PDF物理2頁。

\textbf{公式正答：} 1（single）。

\paragraph{問題の要旨} 熱伝導、計測、熱起電力の記述から最も不適切なものを選ぶ。

\Needspace{4\baselineskip}\paragraph{確認した条件・記述（独自要約）}\mbox{}\par

\noindent\textbf{方向：} 最も不適切。\par

\paragraph{問われる核心} 熱輸送を利用した圧力計測と熱電効果

\paragraph{必要知識} ピラニ計のフィラメント、フォノン、電子、熱電能、フーリエの法則。

\paragraph{考え方・正答を選ぶ理由（査読待ち初稿）} 1は計測の対象と手段を取り違えている。ピラニ計は気体の熱移動による加熱線の熱損失を利用して真空の圧力を測る。2〜5は各材料・線形応答の条件で妥当であり、吸着による圧力変化から熱伝導率を求めるという1を選ぶ。

\Needspace{4\baselineskip}\paragraph{選択肢の検討（査読待ち初稿）}\mbox{}\par

\noindent\textbf{1：誤り} ピラニ真空計は吸着による気体圧力の変化で熱伝導率を求める。ピラニ計は加熱フィラメントから気体へ逃げる熱の圧力依存性を、抵抗や加熱電力で読み圧力へ換算する。吸着を測って熱伝導率を求める説明ではない。\par

\noindent\textbf{2：正しい} 固体の熱伝導はフォノンによるエネルギー移動を伴う。結晶格子の振動励起がエネルギーを運ぶ。金属では伝導電子の寄与もあるので、唯一の担体とは考えない。\par

\noindent\textbf{3：通常の室温材料で正しい} Alの熱伝導率はFeより大きい。産総研の室温付近のデータではAlは226〜237、Feは72〜80.4W/(m・K)程度であり大小が合う。\par

\noindent\textbf{4：正しい} 異種金属の接合部の温度差は導線方向の電場を生む熱起電力に関係する。ゼーベック効果である。異種材料の温度依存する熱電能の差を利用する。\par

\noindent\textbf{5：線形等方近似で正しい} フーリエの法則では熱流密度が温度勾配の大きさに比例する。熱は温度が下がる方向へ流れるのでベクトルとしては温度勾配と逆向きである。\par

\paragraph{再利用する考え方} モデルの前提を明示し、熱輸送を利用した圧力計測と熱電効果から記述や数値を判定する。

\paragraph{共通理論} \hyperref[theory:thermo:transport]{第\ref*{theory:thermo:transport}節}, \hyperref[theory:thermo:measurement]{第\ref*{theory:thermo:measurement}節}

\par 参考：\cite{ulvac_pirani,aist_metal_thermal}。

\par 難易度 1（校準初稿）、優先度未評価。

\subsection{令和元年度再試験 問4}
\label{q:r01-ret-applied-04}

\noindent\textbf{原典：} sources/R01\_retest/17\_applied\_science.pdf、PDF物理2頁。

\textbf{公式正答：} 4（single）。

\paragraph{問題の要旨} 第一法則・第二法則と気体の体積変化を判定し、最も不適切な記述を選ぶ。

\Needspace{4\baselineskip}\paragraph{確認した条件・記述（独自要約）}\mbox{}\par

\noindent\textbf{方向：} 最も不適切。\par

\noindent\textbf{原典の範囲：} 一定質量の気体。理想気体という語は明記されていない。\par

\paragraph{問われる核心} 気体法則とエネルギー保存の対応

\paragraph{必要知識} 第一法則、第二法則、絶対温度と理想気体近似。

\paragraph{考え方・正答を選ぶ理由（査読待ち初稿）} 1と5は第一法則、2は第二法則、3は理想気体近似の定圧関係である。4の等温体積が圧力に比例するという主張は\(V=nRT/p\)に反するので選ぶ。類似問題と肢の位置が異なるため、以前の番号を転用しない。

\Needspace{4\baselineskip}\paragraph{選択肢の検討（査読待ち初稿）}\mbox{}\par

\noindent\textbf{1：正しい} 外部からの熱と系へされた仕事の和が内部エネルギー増加に等しい。熱を流入正、仕事を系へ入力正とすれば\(\Delta U=Q+W_{\mathrm{on}}\)である。系全体の運動・位置エネルギー変化はこの閉じた熱力学系では考えない。\par

\noindent\textbf{2：正しい} 熱は高温から低温へ移り、自然に低温から高温へは移らない。第二法則の巨視的な受動的熱移動の方向である。冷凍機は外部仕事を使うので反例ではない。\par

\noindent\textbf{3：理想気体近似で正しい} 圧力一定・一定質量の気体の体積は絶対温度に比例する。理想気体又はその近似が成り立つ範囲では固定\(n,p\)により\(V=(nR/p)T\)である。\par

\noindent\textbf{4：誤り} 温度一定・一定質量の気体の体積は圧力に比例する。理想気体なら\(pV=nRT\)であり、固定\(n,T\)では\(V=nRT/p\)と反比例する。比例という記述が誤りである。\par

\noindent\textbf{5：正しい} エネルギー変換で関与する全エネルギーの和が保存される。第一法則のエネルギー保存の内容である。対象系から外へ出るエネルギーを数え落とさない。\par

\paragraph{再利用する考え方} モデルの前提を明示し、気体法則とエネルギー保存の対応から記述や数値を判定する。

\paragraph{共通理論} \hyperref[theory:thermo:firstlaw]{第\ref*{theory:thermo:firstlaw}節}, \hyperref[theory:thermo:ideal]{第\ref*{theory:thermo:ideal}節}, \hyperref[theory:thermo:entropy]{第\ref*{theory:thermo:entropy}節}

\par 参考：\cite{mit_thermo_firstlaw2008}。

\paragraph{関連問題} \hyperref[q:h26-reg-applied-08]{H26-REG-APPLIED-08}

\par 難易度 1（校準初稿）、優先度未評価。

\subsection{令和2年度 問6}
\label{q:r02-reg-applied-06}

\noindent\textbf{原典：} sources/R02/17\_applied\_science.pdf、PDF物理2頁。

\textbf{公式正答：} 4（single）。

\paragraph{問題の要旨} 理想気体の定圧・定積4段階変化のうち、第3段階で外部からされる仕事を求める。

\Needspace{4\baselineskip}\paragraph{確認した条件・記述（独自要約）}\mbox{}\par

\noindent\textbf{初期：} \(p_0=2.0\times10^5\,\mathrm{Pa},V_0=100\,\mathrm{m}^3,T_0=27\,{}^\circ\mathrm{C}\)。\par

\noindent\textbf{第1段階：} 定圧で327°Cへ加熱。\par

\noindent\textbf{第2段階：} 定積で27°Cへ冷却。\par

\noindent\textbf{第3段階：} 定圧で−123°Cへ冷却。\par

\noindent\textbf{第4段階：} 定積で27°Cへ加熱。\par

\noindent\textbf{求める量：} 第3段階で気体へされる仕事。\par

\paragraph{問われる核心} 定圧仕事の経路依存と絶対温度の換算

\paragraph{必要知識} 一定気体量のpV/T保存、定積p∝T、定圧V∝T、圧縮仕事の正符号。

\paragraph{考え方・正答を選ぶ理由（査読待ち初稿）} 温度を約300,600,300,150Kとする。第1段階で体積は200m\textsuperscript{3}、第2段階で圧力は\(1.0\times10^5\,\mathrm{Pa}\)になる。第3段階は150Kへの定圧冷却なので体積は100m\textsuperscript{3}となる。気体へされる仕事は\(W_{\mathrm{on}}=p(V_{\mathrm{start}}-V_{\mathrm{end}})=10^5(200-100)=10^7\,\mathrm{J}\)、第4候補である。体積は100m\textsuperscript{3}であり100mLではない。

\Needspace{4\baselineskip}\paragraph{選択肢の検討（査読待ち初稿）}\mbox{}\par

\noindent\textbf{1：条件に不一致} \(1\times10^6\,\mathrm{J}\)。圧縮仕事の計算値より1桁小さい。\par

\noindent\textbf{2：条件に不一致} \(2\times10^6\,\mathrm{J}\)。圧縮の温度幅と気体量の両方を反映していない。\par

\noindent\textbf{3：条件に不一致} \(5\times10^6\,\mathrm{J}\)。第3段階の体積差100m\textsuperscript{3}と圧力10\textasciicircum{}5Paの積の半分である。\par

\noindent\textbf{4：条件に一致} \(1\times10^7\,\mathrm{J}\)。圧力10\textasciicircum{}5Pa×減少体積100m\textsuperscript{3}に一致する。\par

\noindent\textbf{5：条件に不一致} \(2\times10^7\,\mathrm{J}\)。第3段階へ初期圧力2×10\textasciicircum{}5Paをそのまま使った値である。\par

\paragraph{再利用する考え方} モデルの前提を明示し、定圧仕事の経路依存と絶対温度の換算から記述や数値を判定する。

\paragraph{共通理論} \hyperref[theory:thermo:workpaths]{第\ref*{theory:thermo:workpaths}節}, \hyperref[theory:thermo:ideal]{第\ref*{theory:thermo:ideal}節}

\paragraph{関連問題} \hyperref[q:r07-reg-applied-06]{R07-REG-APPLIED-06}

\par 難易度 2（校準初稿）、優先度未評価。

\subsection{令和3年度 問7}
\label{q:r03-reg-applied-07}

\noindent\textbf{原典：} sources/R03/17\_applied\_science.pdf、PDF物理3頁。

\textbf{公式正答：} 1（single）。

\par\textbf{検討未確定：} 1の反比例は明確な誤りで独立選択1。2の室温以上で3Rという一律範囲は高デバイ温度物質で成立しない。3の低温は回転励起・振動凍結の中間範囲、5の100Kも材料依存である。原典の限定不足をAstraの解釈査読対象とし、部分校正を保持する。

\paragraph{問題の要旨} 電子・格子・気体の比熱と相転移の5記述から最も不適切なものを選ぶ。

\Needspace{4\baselineskip}\paragraph{確認した条件・記述（独自要約）}\mbox{}\par

\noindent\textbf{方向：} 最も不適切。\par

\noindent\textbf{記号：} Rは気体定数。\par

\noindent\textbf{原典の省略：} 格子比熱は「室温又はそれ以上」、二原子分子は「比較的低い温度」、固体の熱容量差は「100K以下」と記す。物質や固有温度は指定しない。\par

\paragraph{問われる核心} 量子励起の温度尺度と比熱

\paragraph{必要知識} 電子比熱γT、デバイ温度、二原子分子の自由度、相転移、熱膨張と圧縮率。

\paragraph{考え方・正答を選ぶ理由（査読待ち初稿）} 状態密度\(D(E_F)\)と温度Tへの比例が電子比熱の特徴であり、反比例とする1を最も明確な誤りとして選ぶ。2は高温の古典極限、3は回転のみ励起、4は転移の比熱異常、5は小さな熱膨張に対応する。原典の温度範囲は物質固有の励起温度を記さないため、2・3・5を全材料・全低温に成り立つ法則として説明しない。

\Needspace{4\baselineskip}\paragraph{選択肢の検討（査読待ち初稿）}\mbox{}\par

\noindent\textbf{1：誤り} 電子比熱はフェルミ準位の状態密度と温度に反比例する。通常の縮退電子では\(C_e=(\pi^2/3)k_B^2D(E_F)T\)で、状態密度と温度へ比例し、反比例ではない。\par

\noindent\textbf{2：デバイ温度に比べ高温なら近似的に正しい／原典の範囲に疑義} 格子モル比熱は室温又はそれ以上の高温で3Rとなる。デュロン＝プティ則は\(T\gg\Theta_D\)の古典極限である。室温というだけでは高いデバイ温度を持つ物質には適用できない。\par

\noindent\textbf{3：回転励起・振動凍結の範囲で正しい} 二原子分子の定圧モル比熱は比較的低温で定積の7/5倍である。並進3・回転2自由度なら\(C_{V,m}=5R/2\)、\(C_{p,m}=7R/2\)で比は7/5である。回転まで凍結する極低温や振動の励起される高温では違う。\par

\noindent\textbf{4：二次相転移の比熱異常として妥当} 二次相転移では不連続な変化を伴う異常比熱が現れる。ギブズエネルギーの二階微分に異常が出る。有限跳びや発散があり、全ての連続相転移が同じ形の跳びを持つとは限らない。\par

\noindent\textbf{5：通常の固体での低温近似／物質依存} 固体の定積・定圧モル比熱の差は100K以下で無視できるほど小さい。\(C_p-C_V=TV\alpha_V^2/\kappa_T\)であり低温では一般に小さくなる。ただし100Kという数値だけから全物質で無視可能と断定はできない。\par

\paragraph{再利用する考え方} モデルの前提を明示し、量子励起の温度尺度と比熱から記述や数値を判定する。

\paragraph{共通理論} \hyperref[theory:thermo:degrees]{第\ref*{theory:thermo:degrees}節}, \hyperref[theory:thermo:solidcapacity]{第\ref*{theory:thermo:solidcapacity}節}, \hyperref[theory:thermo:phase]{第\ref*{theory:thermo:phase}節}

\par 参考：\cite{mit_electronic_heat2012,mit_debye2005}。

\par 難易度 2（校準初稿）、優先度未評価。

\subsection{令和4年度 問17}
\label{q:r04-reg-applied-17}

\noindent\textbf{原典：} sources/R04/17\_applied\_science.pdf、PDF物理6頁。

\textbf{公式正答：} 5（single）。

\paragraph{問題の要旨} 実在気体の圧縮係数と分子間力の記述から最も不適切なものを選ぶ。

\Needspace{4\baselineskip}\paragraph{確認した条件・記述（独自要約）}\mbox{}\par

\noindent\textbf{方向：} 最も不適切。\par

\noindent\textbf{定義：} 同温・同圧で実在気体のモル体積を理想気体のモル体積で割った比を圧縮係数とする。\par

\paragraph{問われる核心} 圧縮係数による理想気体からのずれの表現

\paragraph{必要知識} モル体積、理想気体極限、引力と排除体積、ファンデルワールス式。

\paragraph{考え方・正答を選ぶ理由（査読待ち初稿）} 圧縮係数を\(Z=pV_m/(RT)\)と定義する。引力は同温同密度で圧力を下げ、分子の大きさは圧力を上げる方向へ効く。従ってZは1の上下いずれにもなり得て、温度・圧力にかかわらず1より大きいという5は誤りである。

\Needspace{4\baselineskip}\paragraph{選択肢の検討（査読待ち初稿）}\mbox{}\par

\noindent\textbf{1：一般に正しい} 高圧では理想気体の法則から外れる。分子の占有体積や相互作用が無視できず、低密度の理想近似から外れやすい。\par

\noindent\textbf{2：低密度極限で正しい} 低圧では分子間力の影響が小さい。密度が下がると分子間相互作用による圧力への補正は小さくなる。\par

\noindent\textbf{3：正しい} 圧縮係数は同温同圧の実在／理想のモル体積比である。理想モル体積は\(RT/p\)なので比は\(Z=pV_m/(RT)\)である。\par

\noindent\textbf{4：正しい} ファンデルワールス式は引力と分子の大きさを考慮する。\((p+a/V_m^2)(V_m-b)=RT\)でaが引力、bが排除体積の補正を表す。\par

\noindent\textbf{5：誤り} 圧縮係数は温度や圧力に関係なく1より大きい。引力が優勢なら\(Z<1\)、排除体積が優勢なら\(Z>1\)になり得る。低密度では1へ近づく。\par

\paragraph{再利用する考え方} モデルの前提を明示し、圧縮係数による理想気体からのずれの表現から記述や数値を判定する。

\paragraph{共通理論} \hyperref[theory:thermo:realgas]{第\ref*{theory:thermo:realgas}節}

\par 難易度 1（校準初稿）、優先度未評価。

\subsection{令和4年度 問6}
\label{q:r04-reg-applied-06}

\noindent\textbf{原典：} sources/R04/17\_applied\_science.pdf、PDF物理2頁。

\textbf{公式正答：} 1（single）。

\paragraph{問題の要旨} カルノーサイクルの吸熱と2熱源の温度から、1サイクルの低温熱源への放熱を求める。

\Needspace{4\baselineskip}\paragraph{確認した条件・記述（独自要約）}\mbox{}\par

\noindent\textbf{高温熱源：} 1000K\par

\noindent\textbf{低温熱源：} 800K\par

\noindent\textbf{高温側吸熱：} 500J\par

\noindent\textbf{求める量：} 低温源へ放出する熱量の正の大きさ。\par

\paragraph{問われる核心} カルノーサイクルの可逆エントロピー収支

\paragraph{必要知識} 絶対温度、吸熱と放熱の正の大きさ、仕事と熱収支。

\paragraph{考え方・正答を選ぶ理由（査読待ち初稿）} 可逆サイクルのエントロピー変化は0なので\(Q_L=Q_HT_L/T_H\)である。数値を代入すると\(500\times800/1000=400\,\mathrm{J}\)となり、1を選ぶ。取り出した仕事は吸熱と放熱の差であり、求める放熱と混同しない。

\Needspace{4\baselineskip}\paragraph{選択肢の検討（査読待ち初稿）}\mbox{}\par

\noindent\textbf{1：条件に一致} 400J。吸熱に低温／高温の温度比を掛けた放熱に一致する。\par

\noindent\textbf{2：条件に不一致} 450J。可逆性から得る低温側放熱400Jと一致しない。\par

\noindent\textbf{3：条件に不一致} 500J。可逆性から得る低温側放熱400Jと一致しない。\par

\noindent\textbf{4：条件に不一致} 550J。可逆性から得る低温側放熱400Jと一致しない。\par

\noindent\textbf{5：条件に不一致} 600J。可逆性から得る低温側放熱400Jと一致しない。\par

\paragraph{再利用する考え方} モデルの前提を明示し、カルノーサイクルの可逆エントロピー収支から記述や数値を判定する。

\paragraph{共通理論} \hyperref[theory:thermo:carnot]{第\ref*{theory:thermo:carnot}節}

\paragraph{関連問題} \hyperref[q:h27-reg-applied-08]{H27-REG-APPLIED-08}

\par 難易度 2（校準初稿）、優先度未評価。

\subsection{令和4年度 問7}
\label{q:r04-reg-applied-07}

\noindent\textbf{原典：} sources/R04/17\_applied\_science.pdf、PDF物理2頁。

\textbf{公式正答：} 3（single）。

\paragraph{問題の要旨} 熱電対、金属の伝導、多孔質、放射について最も不適切な記述を選ぶ。

\Needspace{4\baselineskip}\paragraph{確認した条件・記述（独自要約）}\mbox{}\par

\noindent\textbf{方向：} 最も不適切。\par

\paragraph{問われる核心} 熱伝導と電気伝導の関係の適用範囲

\paragraph{必要知識} 熱電対、ウィーデマン＝フランツ則、フォノン、多孔質の有効伝導、電磁放射。

\paragraph{考え方・正答を選ぶ理由（査読待ち初稿）} 典型的なバルクダイヤモンドは室温で銀より熱伝導率が大きいので3が不適切である。他の肢は熱電対、金属の電子輸送の近似、多孔質断熱、熱放射に対応する。熱伝導の担体は電子だけではなく、ダイヤモンドでは格子振動が重要である。

\Needspace{4\baselineskip}\paragraph{選択肢の検討（査読待ち初稿）}\mbox{}\par

\noindent\textbf{1：正しい} 異種金属又は合金2種で温度計測できる。熱電対は2種の材料の熱電能の差から温度差を電圧へ変換する。\par

\noindent\textbf{2：通常の金属の電子輸送近似で正しい} 金属の熱伝導率／電気伝導率は極低温を除き絶対温度へ比例する。ウィーデマン＝フランツ則の電子輸送の関係である。散乱の機構や格子熱伝導の大きさによって厳密には変わる。\par

\noindent\textbf{3：通常のバルク比較で誤り} ダイヤモンドの熱伝導率は銀より小さい。高品質バルクダイヤモンドは銀より大きい。金属の電気伝導の序列から熱伝導の序列を全材料へ拡張できない。\par

\noindent\textbf{4：一般的な断熱材料で正しい} 空孔にある気体は熱を運びにくく、多孔質材料は熱伝導率が低い。固体の連続した伝導経路を減らし気体の低い伝導率を利用する。ただし大きな孔の対流や高温の放射は別に評価する。\par

\noindent\textbf{5：正しい} 熱放射は電磁波でのエネルギー移動である。媒体の巨視的移動が必要な対流とは区別する。\par

\paragraph{再利用する考え方} モデルの前提を明示し、熱伝導と電気伝導の関係の適用範囲から記述や数値を判定する。

\paragraph{共通理論} \hyperref[theory:thermo:transport]{第\ref*{theory:thermo:transport}節}, \hyperref[theory:thermo:measurement]{第\ref*{theory:thermo:measurement}節}

\par 参考：\cite{aist_metal_thermal,wang_sun_diamond2024}。

\par 難易度 1（校準初稿）、優先度未評価。

\subsection{令和5年度 問6}
\label{q:r05-reg-applied-06}

\noindent\textbf{原典：} sources/R05/17\_applied\_science.pdf、PDF物理3頁。

\textbf{公式正答：} 3（single）。

\paragraph{問題の要旨} 相転移の次数、熱力学量、強磁性、過冷却の記述から最も不適切なものを選ぶ。

\Needspace{4\baselineskip}\paragraph{確認した条件・記述（独自要約）}\mbox{}\par

\noindent\textbf{方向：} 最も不適切。\par

\noindent\textbf{対象：} 融解・気化と連続相転移の通常の分類。\par

\paragraph{問われる核心} 自由エネルギーの微分から見る相転移と準安定状態

\paragraph{必要知識} ギブズエネルギー、潜熱、エントロピー、熱容量、秩序変数、核生成。

\paragraph{考え方・正答を選ぶ理由（査読待ち初稿）} 相転移の次数は自由エネルギーの微分の性質で整理する。一次ではエントロピーなど一階微分が跳び、二次・連続転移では比熱など二階微分が異常を示す。熱容量は連続であると一律に述べる3が不適切である。過冷却は平衡相図と変化の速度を区別する例である。

\Needspace{4\baselineskip}\paragraph{選択肢の検討（査読待ち初稿）}\mbox{}\par

\noindent\textbf{1：臨界点から離れた相共存で正しい} 融解と気化は一次相転移である。融解・気化は潜熱を伴う通常の一次転移である。液気臨界点で相の区別がなくなる場合は別である。\par

\noindent\textbf{2：正しい} 一次相転移ではエントロピーが不連続に変わる。潜熱Lと転移温度Tにより\(\Delta S=L/T\)でエントロピーの跳びがある。\par

\noindent\textbf{3：誤り} 二次相転移では熱容量は連続に変わる。二次転移ではギブズエネルギーの一階微分は連続でも、比熱のような二階微分に跳びや発散が出る。\par

\noindent\textbf{4：通常の零磁場のキュリー転移として正しい} 鉄の強磁性から常磁性への転移は二次相転移である。磁化という秩序変数が零磁場で連続に消える転移であり、比熱には異常が現れる。\par

\noindent\textbf{5：正しい} 液体や気体を平衡転移温度より冷やしても相変化しない過冷却がある。新相の核生成の障壁などで準安定状態が続く。平衡の共存温度そのものが変わったとの意味ではない。\par

\paragraph{再利用する考え方} モデルの前提を明示し、自由エネルギーの微分から見る相転移と準安定状態から記述や数値を判定する。

\paragraph{共通理論} \hyperref[theory:thermo:phase]{第\ref*{theory:thermo:phase}節}

\par 難易度 1（校準初稿）、優先度未評価。

\subsection{令和6年度 問6}
\label{q:r06-reg-applied-06}

\noindent\textbf{原典：} sources/R06/17\_applied\_science.pdf、PDF物理2頁。

\textbf{公式正答：} 3（single）。

\paragraph{問題の要旨} ヒーターで1molの物質Aを一定外圧で全て気体にしたとき、物質Aが外部へ行う仕事を求める。

\Needspace{4\baselineskip}\paragraph{確認した条件・記述（独自要約）}\mbox{}\par

\noindent\textbf{系：} 断熱シリンダー内の1molの液体物質Aと、加熱用ヒーター。\par

\noindent\textbf{温度：} 物質Aの沸点506K\par

\noindent\textbf{外圧：} 0.1MPa一定\par

\noindent\textbf{仮定：} 液体体積を無視、気体状態の物質Aは理想気体。\par

\noindent\textbf{気体定数：} 8.3J/(mol・K)\par

\noindent\textbf{求める量：} 物質Aが外部へ行う仕事の正の大きさ。\par

\paragraph{問われる核心} 気化で増えた体積による定圧仕事

\paragraph{必要知識} 一定外圧の仕事、理想気体、液体体積無視、壁からの熱流と電気入力の区別。

\paragraph{考え方・正答を選ぶ理由（査読待ち初稿）} 外部への体積仕事は\(W_{\mathrm{out}}=p_{\mathrm{ext}}(V_g-V_l)\)である。液体体積を無視し、終状態の理想気体に\(p_{\mathrm{ext}}V_g=nRT\)を適用すると\(W_{\mathrm{out}}=1\times8.3\times506=4199.8\,\mathrm{J}\)である。最も近い4200Jの3を選ぶ。加熱の総エネルギーや潜熱を問う問題ではない。断熱壁でもヒーターから仕事としてエネルギーが入る。

\Needspace{4\baselineskip}\paragraph{選択肢の検討（査読待ち初稿）}\mbox{}\par

\noindent\textbf{1：条件に不一致} 830J。外圧×体積増加=nRTに対して桁又は値が合わない。\par

\noindent\textbf{2：条件に不一致} 2240J。外圧×体積増加=nRTに対して桁又は値が合わない。\par

\noindent\textbf{3：条件に一致} 4200J。理想蒸気のpV=\allowbreak{}nRTに一致する。\par

\noindent\textbf{4：条件に不一致} 8300J。外圧×体積増加=nRTに対して桁又は値が合わない。\par

\noindent\textbf{5：条件に不一致} 22400J。外圧×体積増加=nRTに対して桁又は値が合わない。\par

\paragraph{再利用する考え方} モデルの前提を明示し、気化で増えた体積による定圧仕事から記述や数値を判定する。

\paragraph{共通理論} \hyperref[theory:thermo:firstlaw]{第\ref*{theory:thermo:firstlaw}節}, \hyperref[theory:thermo:workpaths]{第\ref*{theory:thermo:workpaths}節}

\par 参考：\cite{mit_thermo_firstlaw2008}。

\paragraph{関連問題} \hyperref[q:h29-reg-applied-06]{H29-REG-APPLIED-06}

\par 難易度 2（校準初稿）、優先度未評価。

\subsection{令和7年度 問6}
\label{q:r07-reg-applied-06}

\noindent\textbf{原典：} sources/R07/17\_applied\_science.pdf、PDF物理3頁。

\textbf{公式正答：} 1（single）。

\paragraph{問題の要旨} 理想気体を定圧・定積で4段階に変化させ、第3段階の外部からの仕事を求める。

\Needspace{4\baselineskip}\paragraph{確認した条件・記述（独自要約）}\mbox{}\par

\noindent\textbf{初期：} \(p_0=5.0\times10^3\,\mathrm{Pa},V_0=120\,\mathrm{m}^3,T_0=27\,{}^\circ\mathrm{C}\)。\par

\noindent\textbf{第1段階：} 定圧で227°Cへ加熱。\par

\noindent\textbf{第2段階：} 定積で27°Cへ冷却。\par

\noindent\textbf{第3段階：} 定圧で−73°Cへ冷却。\par

\noindent\textbf{第4段階：} 定積で27°Cへ加熱。\par

\noindent\textbf{求める量：} 第3段階で気体へされる仕事。\par

\paragraph{問われる核心} 状態方程式から気体量を消去する定圧仕事

\paragraph{必要知識} pV=\allowbreak{}nRT、絶対温度、定圧仕事−nRΔT、仕事の向き。

\paragraph{考え方・正答を選ぶ理由（査読待ち初稿）} 絶対温度を約300,500,300,200Kとする。気体量を\(nR=p_0V_0/T_0=(5.0\times10^3)120/300=2000\,\mathrm{J/K}\)と置く。第3段階の定圧圧縮仕事は\(W_{\mathrm{on}}=-nR\Delta T=2000(300-200)=2.0\times10^5\,\mathrm{J}\)、第1候補である。途中状態はV=\allowbreak{}200m\textsuperscript{3}、p=\allowbreak{}3000Pa、終体積400/3m\textsuperscript{3}であり、pΔVからも同じ値になる。

\Needspace{4\baselineskip}\paragraph{選択肢の検討（査読待ち初稿）}\mbox{}\par

\noindent\textbf{1：条件に一致} \(2\times10^5\,\mathrm{J}\)。nR=初期pV/Tと圧縮温度幅100Kの積に一致する。\par

\noindent\textbf{2：条件に不一致} \(5\times10^5\,\mathrm{J}\)。原典の圧力の指数と圧縮温度幅を用いた結果と合わない。\par

\noindent\textbf{3：条件に不一致} \(1\times10^6\,\mathrm{J}\)。計算値の5倍であり、段階の温度と圧力を区別していない。\par

\noindent\textbf{4：条件に不一致} \(2\times10^6\,\mathrm{J}\)。計算値より1桁大きい。初期圧力の指数は3である。\par

\noindent\textbf{5：条件に不一致} \(5\times10^6\,\mathrm{J}\)。圧力の指数又は体積変化の計算が一致しない。\par

\paragraph{再利用する考え方} モデルの前提を明示し、状態方程式から気体量を消去する定圧仕事から記述や数値を判定する。

\paragraph{共通理論} \hyperref[theory:thermo:workpaths]{第\ref*{theory:thermo:workpaths}節}, \hyperref[theory:thermo:ideal]{第\ref*{theory:thermo:ideal}節}

\paragraph{関連問題} \hyperref[q:r02-reg-applied-06]{R02-REG-APPLIED-06}

\par 難易度 2（校準初稿）、優先度未評価。

\section{相平衡・溶液}

\subsection{平成25年度 問9}
\label{q:h25-reg-applied-09}

\noindent\textbf{原典：} sources/H25/17\_applied\_science.pdf、PDF物理5頁。

\textbf{公式正答：} 4（single）。

\par\textbf{検討未確定：} 原典の断熱膨張の肢は外部仕事又は可逆性を明記しない。理想気体の断熱自由膨張は温度不変であり、この肢も無条件には正しくない。熱容量の逆転肢を最も明確な誤りとして独立選択したが、原典の省略を追加条件で書き換えず部分校正としてAstraへ送る。

\paragraph{問題の要旨} 理想気体と溶液・相平衡の記述から、最も不適切なものを選ぶ。

\Needspace{4\baselineskip}\paragraph{確認した条件・記述（独自要約）}\mbox{}\par

\noindent\textbf{方向：} 最も不適切。\par

\noindent\textbf{前提の確認：} 断熱膨張の肢に外部仕事又は可逆性の指定はない。\par

\paragraph{問われる核心} 熱容量と膨張仕事・相平衡の成立条件

\paragraph{必要知識} 理想気体の熱容量差、断熱と自由膨張、液気平衡、混合エントロピー、ラウールの法則。

\paragraph{考え方・正答を選ぶ理由（査読待ち初稿）} 同じ理想気体を同じだけ昇温させると定積では\(Q_V=nC_{V,m}\Delta T\)、定圧では\(Q_p=nC_{p,m}\Delta T\)を要し、\(C_{p,m}-C_{V,m}=R>0\)である。定積の方が熱量が大きいという肢4を選ぶ。断熱膨張の冷却は正の外部仕事を行う場合に限り、自由膨張という反例があるので解釈疑義を保持する。

\Needspace{4\baselineskip}\paragraph{選択肢の検討（査読待ち初稿）}\mbox{}\par

\noindent\textbf{1：通常の液気平衡で正しい} 高圧ほど純液体の沸点が上がる。液気平衡で\(dp/dT=L/(T\Delta V)>0\)なので外圧上昇で沸点が上がる。臨界点から離れた通常の沸騰を考える。\par

\noindent\textbf{2：仕事を外へ行う膨張なら正しい／前提省略} 理想気体を断熱膨張させると温度が下がる。断熱では\(\Delta U=-W_{\mathrm{out}}\)。正の仕事を外へ行う理想気体は冷えるが、真空への自由膨張なら\(W=0\)で温度は変わらない。膨張だけで冷却は保証されない。\par

\noindent\textbf{3：理想的な異種気体で正しい} 同温・同圧で異種気体の混合が自発的に起こる。非反応の異種理想気体の混合では\(\Delta G=nRT\sum_i x_i\ln x_i<0\)となる。\par

\noindent\textbf{4：誤り} 理想気体の同じ温度上昇には定積の方が定圧より多い熱が必要。定圧では膨張仕事も必要であり、理想気体のモル熱容量は定圧が定積よりR大きい。定積の方が多いという記述は逆である。\par

\noindent\textbf{5：理想溶液で正しい} 理想溶液の各成分の蒸気圧は純成分の蒸気圧×液相モル分率である。ラウールの法則\(p_i=x_i p_i^*\)である。実在溶液の全組成範囲へ無条件に使う法則ではない。\par

\paragraph{再利用する考え方} モデルの前提を明示し、熱容量と膨張仕事・相平衡の成立条件から記述や数値を判定する。

\paragraph{共通理論} \hyperref[theory:thermo:capacity]{第\ref*{theory:thermo:capacity}節}, \hyperref[theory:thermo:adiabatic]{第\ref*{theory:thermo:adiabatic}節}, \hyperref[theory:thermo:phase]{第\ref*{theory:thermo:phase}節}, \hyperref[theory:thermo:mixing]{第\ref*{theory:thermo:mixing}節}

\par 参考：\cite{mit_thermo_adiabatic2008,mit_clapeyron2008,mit_solutions2008}。

\paragraph{関連問題} \hyperref[q:h30-reg-applied-06]{H30-REG-APPLIED-06}

\par 難易度 2（校準初稿）、優先度未評価。

\subsection{平成30年度 問6}
\label{q:h30-reg-applied-06}

\noindent\textbf{原典：} sources/H30/17\_applied\_science.pdf、PDF物理2頁。

\textbf{公式正答：} 5（single）。

\par\textbf{検討未確定：} 原典の断熱膨張の肢は外部仕事又は可逆性を明記しない。理想気体の断熱自由膨張は温度不変であり、この肢も無条件には正しくない。熱容量の逆転肢を最も明確な誤りとして独立選択したが、原典の省略を追加条件で書き換えず部分校正としてAstraへ送る。

\paragraph{問題の要旨} 理想気体と溶液・相平衡の記述から、最も不適切なものを選ぶ。

\Needspace{4\baselineskip}\paragraph{確認した条件・記述（独自要約）}\mbox{}\par

\noindent\textbf{方向：} 最も不適切。\par

\noindent\textbf{前提の確認：} 断熱膨張の肢に外部仕事又は可逆性の指定はない。\par

\paragraph{問われる核心} 熱容量と膨張仕事・相平衡の成立条件

\paragraph{必要知識} 理想気体の熱容量差、断熱と自由膨張、液気平衡、混合エントロピー、ラウールの法則。

\paragraph{考え方・正答を選ぶ理由（査読待ち初稿）} 同じ理想気体を同じだけ昇温させると定積では\(Q_V=nC_{V,m}\Delta T\)、定圧では\(Q_p=nC_{p,m}\Delta T\)を要し、\(C_{p,m}-C_{V,m}=R>0\)である。定積の方が熱量が大きいという肢5を選ぶ。断熱膨張の冷却は正の外部仕事を行う場合に限り、自由膨張という反例があるので解釈疑義を保持する。

\Needspace{4\baselineskip}\paragraph{選択肢の検討（査読待ち初稿）}\mbox{}\par

\noindent\textbf{1：仕事を外へ行う膨張なら正しい／前提省略} 理想気体を断熱膨張させると温度が下がる。断熱では\(\Delta U=-W_{\mathrm{out}}\)。正の仕事を外へ行う理想気体は冷えるが、真空への自由膨張なら\(W=0\)で温度は変わらない。膨張だけで冷却は保証されない。\par

\noindent\textbf{2：理想的な異種気体で正しい} 同温・同圧で異種気体の混合が自発的に起こる。非反応の異種理想気体の混合では\(\Delta G=nRT\sum_i x_i\ln x_i<0\)となる。\par

\noindent\textbf{3：理想溶液で正しい} 理想溶液の各成分の蒸気圧は純成分の蒸気圧×液相モル分率である。ラウールの法則\(p_i=x_i p_i^*\)である。実在溶液の全組成範囲へ無条件に使う法則ではない。\par

\noindent\textbf{4：通常の液気平衡で正しい} 高圧ほど純液体の沸点が上がる。液気平衡で\(dp/dT=L/(T\Delta V)>0\)なので外圧上昇で沸点が上がる。臨界点から離れた通常の沸騰を考える。\par

\noindent\textbf{5：誤り} 理想気体の同じ温度上昇には定積の方が定圧より多い熱が必要。定圧では膨張仕事も必要であり、理想気体のモル熱容量は定圧が定積よりR大きい。定積の方が多いという記述は逆である。\par

\paragraph{再利用する考え方} モデルの前提を明示し、熱容量と膨張仕事・相平衡の成立条件から記述や数値を判定する。

\paragraph{共通理論} \hyperref[theory:thermo:capacity]{第\ref*{theory:thermo:capacity}節}, \hyperref[theory:thermo:adiabatic]{第\ref*{theory:thermo:adiabatic}節}, \hyperref[theory:thermo:phase]{第\ref*{theory:thermo:phase}節}, \hyperref[theory:thermo:mixing]{第\ref*{theory:thermo:mixing}節}

\par 参考：\cite{mit_thermo_adiabatic2008,mit_clapeyron2008,mit_solutions2008}。

\paragraph{関連問題} \hyperref[q:h25-reg-applied-09]{H25-REG-APPLIED-09}

\par 難易度 2（校準初稿）、優先度未評価。

\subsection{令和4年度 問11}
\label{q:r04-reg-applied-11}

\noindent\textbf{原典：} sources/R04/17\_applied\_science.pdf、PDF物理4頁。

\textbf{公式正答：} 2（single）。

\paragraph{問題の要旨} 水の三重点を既知として、富士山頂の大気圧下の沸点を求めるための理論式を選ぶ。

\Needspace{4\baselineskip}\paragraph{確認した条件・記述（独自要約）}\mbox{}\par

\noindent\textbf{既知のもの：} 水の三重点の温度と蒸気圧、山頂の大気圧。\par

\noindent\textbf{求める対象：} 純水の液気共存温度。数値の沸点を答える問題ではない。\par

\paragraph{問われる核心} 相共存曲線と圧力による沸点の変化

\paragraph{必要知識} 相平衡の化学ポテンシャル一致、潜熱、蒸気圧、理想蒸気近似。

\paragraph{考え方・正答を選ぶ理由（査読待ち初稿）} 沸点では蒸気圧が外圧に等しく、液気共存曲線をたどればよい。関係式は\(dp/dT=L/[T(V_g-V_l)]\)で、蒸気を理想気体とし液体体積を無視すると\(d\ln p/dT=L_m/(RT^2)\)となる。クラウジウス＝クラペイロンの式の2を選ぶ。有限区間の数値を求めるには潜熱やその温度依存なども必要であり、指定されない山頂の沸点を補わない。

\Needspace{4\baselineskip}\paragraph{選択肢の検討（査読待ち初稿）}\mbox{}\par

\noindent\textbf{1：条件に不一致} ボルンの式。イオンの静電的な溶媒和などを扱う式で、直接の液気共存曲線ではない。\par

\noindent\textbf{2：条件に一致} クラウジウス＝クラペイロンの式。液気共存の圧力と温度の関係を与える。\par

\noindent\textbf{3：条件に不一致} ファントホッフの式。平衡定数の温度依存を扱うが、この液気平衡を直接記す名称としてはクラウジウス＝クラペイロンが適切である。\par

\noindent\textbf{4：条件に不一致} プランクの式。黒体放射のスペクトルを扱う。\par

\noindent\textbf{5：条件に不一致} レイリー＝ジーンズの式。黒体放射の古典的な低周波近似を扱う。\par

\paragraph{再利用する考え方} モデルの前提を明示し、相共存曲線と圧力による沸点の変化から記述や数値を判定する。

\paragraph{共通理論} \hyperref[theory:thermo:phase]{第\ref*{theory:thermo:phase}節}

\par 参考：\cite{mit_clapeyron2008}。

\par 難易度 1（校準初稿）、優先度未評価。

